\documentclass[../Main_Include.tex]{subfiles}
\begin{document}
\backgroundsetup{contents={}}  % suprime draft de background en compilación standalone
% ============================
% LABORATORIO 1 — Laboratorio1_1010.tex
% ============================
%
% CONTENIDO:
%   Sec. 1 — Ejercicio 1: Chatbot básico de IA
%   Sec. 2 — Ejercicio 2: Chatbot mejorado con difflib
%
% DEPENDENCIAS (cargadas desde Main_Include.tex):
%   Preambulo.tex       — paquetes, colores, estilos de código
%   PaginaEstilo.tex    — lstset global
%   EstiloSecciones.tex — \titleformat, fancyhdr
%   Entornos.tex        — cuadropregunta
%
% ============================

% --- Título del encabezado para este informe ---
\newcommand{\tituloinforme}{Fundamentos de la Inteligencia Artificial — Laboratorio 1}

% ══════════════════════════════════════════════════════════════════════
\section{Ejercicio 1 — Chatbot básico de IA}

Se implementó un chatbot en Python que responde preguntas básicas sobre IA
mediante un diccionario de pares clave--respuesta. La entrada del usuario
se convierte a minúsculas y se compara directamente con las claves definidas.

% ──────────────────────────────────────────────────────────────────────
\subsection*{Observaciones}

El chatbot responde correctamente únicamente cuando la pregunta se escribe
\textbf{con tilde}, tal como están definidas las claves del diccionario.
Sin tilde la comparación falla y el programa devuelve
\texttt{"No entiendo tu pregunta."}.

Las preguntas válidas son:

\begin{itemize}[leftmargin=1.5cm, itemsep=2pt]
    \item \texttt{qué es la ia} → devuelve la definición de IA.
    \item \texttt{quién inventó la ia} → devuelve la referencia a John McCarthy.
    \item \texttt{adiós} → cierra el programa.
\end{itemize}

% ──────────────────────────────────────────────────────────────────────
\subsection*{Capturas de pantalla}

\begin{figure}[H]
    \centering
    \includegraphics[width=\linewidth]{Graficos/Lab1/E1_fun_001.png}
    \caption{Ejecución del chatbot — se observan los fallos sin tilde y
             las respuestas exitosas con tilde.}
    \label{fig:e1_001}
\end{figure}

% ══════════════════════════════════════════════════════════════════════
\newpage
\section{Ejercicio 2 — Chatbot mejorado con \texttt{difflib}}

Se amplió el chatbot incorporando la librería \texttt{difflib} para tolerancia
ortográfica y una interfaz de widgets interactiva mediante \texttt{ipywidgets},
diseñada para ejecutarse en Google Colab. El bot acepta preguntas con errores
menores de escritura y permite al usuario enseñarle nuevas respuestas en tiempo
de ejecución, que se persisten en un archivo \texttt{respuestas.json}.

% ──────────────────────────────────────────────────────────────────────
\subsection*{Observaciones}

A diferencia del Ejercicio 1, la función \texttt{difflib.get\_close\_matches}
permite reconocer preguntas sin tilde o con ortografía aproximada.
Por ejemplo, \texttt{que es la ia} (sin tilde) devuelve la respuesta correcta.

Las preguntas base del diccionario son:

\begin{itemize}[leftmargin=1.5cm, itemsep=2pt]
    \item \texttt{qué es la ia}
    \item \texttt{quién inventó la ia}
    \item \texttt{qué es machine learning}
    \item \texttt{qué es deep learning}
    \item \texttt{qué diferencia hay entre machine learning y deep learning}
\end{itemize}

\paragraph{Caso borde — entrada vacía.}
Al presionar \emph{Enviar} sin texto, el panel de aprendizaje queda abierto
sin posibilidad de cerrarlo, salvo ingresando y guardando una respuesta válida.

\paragraph{Persistencia.}
Las respuestas aprendidas se guardan en \texttt{respuestas.json} dentro del
entorno de Colab. El archivo es volátil: al cerrar la sesión se elimina y el
chatbot vuelve al diccionario base.

\paragraph{Comando de salida.}
\texttt{adiós} (con tilde) muestra \texttt{"¡Hasta luego!"} y detiene el procesamiento.

% ──────────────────────────────────────────────────────────────────────
\subsection*{Capturas de pantalla}

\begin{figure}[H]
    \centering
    \includegraphics[width=\linewidth]{Graficos/Lab1/E2_fun_001.png}
    \caption{Vista del código en Google Colab — diccionario base y función
             \texttt{guardar\_respuestas}.}
    \label{fig:e2_001}
\end{figure}

\begin{figure}[H]
    \centering
    \includegraphics[width=\linewidth]{Graficos/Lab1/E2_fun_002.png}
    \caption{Tolerancia ortográfica — \texttt{que es la ia} (sin tilde)
             devuelve la respuesta correcta.}
    \label{fig:e2_002}
\end{figure}

\begin{figure}[H]
    \centering
    \includegraphics[width=\linewidth]{Graficos/Lab1/E2_fun_003.png}
    \caption{Segunda prueba de tolerancia — \texttt{quien invento la ia}
             (sin tilde ni acento) reconocida correctamente.}
    \label{fig:e2_003}
\end{figure}

\begin{figure}[H]
    \centering
    \includegraphics[width=\linewidth]{Graficos/Lab1/E2_fun_004.png}
    \caption{Pregunta desconocida — el bot no reconoce \texttt{adios} y despliega
             el panel de aprendizaje para ingresar y guardar una nueva respuesta.}
    \label{fig:e2_004}
\end{figure}

\begin{figure}[H]
    \centering
    \includegraphics[width=\linewidth]{Graficos/Lab1/E2_fun_005.png}
    \caption{Comando de salida — \texttt{adiós} con tilde cierra el procesamiento.}
    \label{fig:e2_005}
\end{figure}

\begin{figure}[H]
    \centering
    \includegraphics[width=\linewidth]{Graficos/Lab1/E2_fun_006.png}
    \caption{Contenido del archivo \texttt{respuestas.json} tras la sesión —
             se verifican las respuestas aprendidas persistidas.}
    \label{fig:e2_006}
\end{figure}

% ──────────────────────────────────────────────────────────────────────
\newpage
\subsection*{Verificación de persistencia}
\vspace{0.5cm}
Tras la sesión se inspeccionó el archivo \texttt{respuestas.json} ejecutando
el siguiente comando en una celda de Colab:
\vspace{0.3cm}
\begin{lstlisting}[style=py_dark, label={lst:json_check}]
import json
with open("respuestas.json", "r", encoding="utf-8") as f:
    print(json.dumps(json.load(f), ensure_ascii=False, indent=4))
\end{lstlisting}

\vspace{0.5cm}
La salida obtenida fue:
\vspace{0.3cm}
\begin{tcolorbox}[
    colback=bg_matlab,
    colframe=azulprinc,
    fontupper=\ttfamily\small,
    boxrule=0.4pt,
    left=6pt, right=6pt, top=4pt, bottom=4pt
]
\begin{verbatim}
{
    "qué es la ia": "La Inteligencia Artificial es la simulación
                     de la inteligencia humana en computadoras.",
    "quién inventó la ia": "El término 'IA' fue acuñado
                            por John McCarthy en 1956.",
    "qué es machine learning": "Machine Learning es un subconjunto
                                de la IA que permite a las máquinas
                                aprender a partir de datos sin ser
                                programadas explícitamente.",
    "qué es deep learning": "Deep Learning es un tipo de Machine
                             Learning basado en redes neuronales
                             profundas.",
    "qué diferencia hay entre machine learning y deep learning":
        "Machine Learning abarca diferentes técnicas de aprendizaje
         automático, mientras que Deep Learning usa redes neuronales
         profundas para procesar datos complejos.",
    "adios": "chao, te cuidas",
    "nos vemos": "bye",
    "¡hasta luego!": "x"
}
\end{verbatim}
\end{tcolorbox}

Se verifica que las respuestas aprendidas durante la sesión ---
\texttt{adios}, \texttt{nos vemos} y \texttt{¡hasta luego!} ---
fueron correctamente persistidas en el archivo JSON.

% ══════════════════════════════════════════════════════════════════════
\newpage
\appendix
\renewcommand{\thesection}{\Alph{section}}
\section*{Apéndice — Códigos fuente}
\addcontentsline{toc}{section}{Apéndice}

\subsection*{Ejercicio 1 — Chatbot básico}

\lstinputlisting[style=py_dark,
    caption={Ejercicio 1 --- Chatbot b\'asico de IA},
    label={lst:e1}]{Anexos/Lab1/E1_FUNDAMENTOS_ETN1010.py}

\newpage

\subsection*{Ejercicio 2 — Chatbot mejorado}

\lstinputlisting[style=py_dark,
    caption={Ejercicio 2 --- Chatbot mejorado con \texttt{difflib} e \texttt{ipywidgets}},
    label={lst:e2}]{Anexos/Lab1/E2_FUNDAMENTOS_ETN1010.py}

\end{document}
